% Upload this single file to Overleaf and choose pdfLaTeX.
% The bibliography is embedded. No .bib file or external images are needed.
% Six substantive sections start on separate pages, followed by references.
% Author and affiliation below come from the supplied project presentation.
% Recorded measurements come from artifacts/pretrained-hatexplain.json.
\documentclass[11pt,a4paper]{article}
\usepackage[T1]{fontenc}
\usepackage[utf8]{inputenc}
\usepackage[margin=1in]{geometry}
\usepackage{mathptmx}
\usepackage{amsmath,amssymb}
\usepackage{booktabs,tabularx,array}
\usepackage{microtype}
\usepackage{xurl}
\usepackage[hidelinks]{hyperref}
\usepackage{enumitem}
\setlist{itemsep=2pt,topsep=4pt}
\setlength{\parskip}{4pt}
\setlength{\parindent}{1em}
\setlength{\emergencystretch}{2em}
\newcolumntype{Y}{>{\raggedright\arraybackslash}X}
\newcommand{\artifact}[1]{\texttt{\detokenize{#1}}}

\title{Hate Speech Detection through Explainable AI and Moderator Feedback\\[4pt]
\large An NLP Prototype and a Framework for Reward Optimization}
\author{Moksha Sharma\\Department of Computer Science and Engineering\\Bennett University}
\date{November 2026}

\begin{document}
\maketitle
\begin{abstract}
Hate speech detection requires decisions that moderators can inspect and correct.
This paper presents an NLP prototype that combines a pretrained text classifier,
word attributions, and persistent moderator feedback. The executable application
uses a downloaded TF-IDF and logistic regression pipeline. SHAP and LIME explain
its class probabilities through token deletion, while a review interface records
confirmations and corrections. An existing evaluation on 1,923 deduplicated
HateXplain test posts gives 62.77\% accuracy, 42.36\% hate precision, 56.90\% hate
recall, and 48.56\% hate F1 at a fixed threshold of 0.50. SHAP achieves mean
token-overlap F1 of 0.513 on a sample of 30 hateful posts with human rationales.
The publisher's precise training corpus and split are unspecified, so these
measurements describe behavior on the evaluation corpus with possible training
overlap. We also specify an offline preference-optimization objective and
describe source modules for a classical SVM baseline, multilingual transformer,
and image--caption classification. Their performance remains to be established.
The prototype provides an inspectable moderation workflow and a concrete
experimental basis for studying classification errors, explanation quality,
and feedback-driven adaptation.
\end{abstract}
\noindent\textbf{Keywords}\quad Natural language processing, hate speech,
explainable AI, SHAP, LIME, human feedback, preference optimization.

\section{Introduction}
Online moderation involves language whose meaning depends on context, target,
and communicative intent. A post can contain offensive vocabulary while quoting
or condemning discrimination. Conversely, an exclusionary statement can target
a community through indirect expressions. Davidson et al. demonstrate why hate
and offensive language require distinct labels \cite{davidson}. This distinction
motivates the binary task used here, where hate is the positive class and other
source labels are combined into the negative class.

A prediction score provides limited assistance when the reviewer cannot inspect
which parts of a post affected it. Explanation methods can expose influential
words, and a review interface can capture a correction with its original model
context. However, visible highlights need evaluation. They can reflect lexical
shortcuts or perturbation artifacts, and agreement with an annotation captures
only one aspect of explanation quality \cite{eraser}.

The project investigates how classification, explanation, and recorded review
can operate in one reproducible workflow. Its measured contribution is an
integrated prototype and an evaluation of a frozen third-party classifier. Its
research framework extends that workflow with offline preference optimization,
multilingual modeling, and image--caption experiments. The evaluation answers
three questions. How does the available classifier behave under a consistent
hate-only label mapping? How closely do supportive word attributions align with
annotated evidence? What information must a review record preserve to support
a later, auditable model update?

The resulting system exposes the predicted label, hate probability, configurable
threshold, signed word effects, and review controls. The paper reports the
available measurements and identifies the additional experiments required for
claims about adaptation, generalization, or multilingual coverage.

\clearpage
\section{Related Work}
\subsection{Hate speech datasets and contextual models}
HateXplain supplies class labels, target annotations, and human rationales,
supporting joint examination of classification and explanatory evidence
\cite{hatexplain}. The Davidson dataset distinguishes hate speech, offensive
language, and neither \cite{davidson}. HASOC expands the dataset landscape to
Indo-European languages and illustrates the importance of edition-specific
task definitions \cite{hasoc}. A shared binary task therefore requires explicit
mapping rules. Combining an offensive category with hate would change the
research question and the meaning of the reported precision.

The Transformer architecture introduced self-attention as a general mechanism
for sequence modeling \cite{attention}. BERT adapts bidirectional pretraining
to downstream language tasks \cite{bert}. HateBERT investigates additional
pretraining on abusive English content, demonstrating the relevance of domain
adaptation for this task family \cite{hatebert}. XLM-R provides a multilingual
encoder that can support later experiments involving Hindi and English
\cite{xlmr}. Multilingual pretraining alone leaves open whether a classifier
has learned the local categories and code-mixed expressions required by a
particular moderation setting.

The Hateful Memes benchmark introduces examples where image and caption
must be considered jointly \cite{memes}. It motivates the project's separate
image--caption route, whose evaluation includes both single-modality models
and a fused model. This comparison is necessary to determine whether a system
uses information from both modalities.

\subsection{Explanations and diagnostic evaluation}
LIME estimates a simple model near an input using weighted perturbations
\cite{lime}. SHAP provides an additive feature-attribution framework grounded
in Shapley values \cite{shap}. Both can query a common probability interface,
making them suitable for comparing classifier backends. Their outputs depend
on the feature representation and on how perturbed inputs are constructed.

ERASER distinguishes agreement with human rationales from faithfulness to a
model's prediction \cite{eraser}. HARE studies step-by-step reasoning for
explainable hate speech detection \cite{hare}. The present application uses
token attributions, while natural-language reasoning remains a possible
extension requiring its own verification. HateCheck provides functional tests
that complement aggregate accuracy and F1 \cite{hatecheck}. Such tests are
relevant for quotation, negation, identity mentions, and implicit attacks.

\subsection{Feedback, fairness, and reporting}
Instruction tuning with human feedback demonstrates a route for adapting
language models to reviewed outputs \cite{ouyang}. Direct Preference
Optimization (DPO) formulates preference learning through probability ratios
relative to a reference model \cite{dpo}. Our proposed adaptation uses a
reviewed binary label as the preferred answer, which requires evaluation on
this classification setting before any improvement claim.

Sap et al. identify racial bias risks in hate speech detection \cite{sap}.
This motivates slice diagnostics and targeted error inspection. Model cards
and dataset datasheets offer documentation frameworks for intended use,
collection conditions, and limitations \cite{modelcards,datasheets}.
The prototype follows these reporting goals by recording model provenance,
label mappings, and evaluation scope.

\clearpage
\section{Data Preparation and Classification Method}
\subsection{Task and canonical representation}
Let $x$ denote a social media post and $y\in\{0,1\}$ its target label. The
project maps hateful content to $1$ and the remaining mapped categories to
$0$. A canonical record contains an identifier, text, label, language,
optional annotated target groups, optional token rationales, and an optional
split assignment. Multiple groups can occur in one post.

The HateXplain adapter uses a strict majority over the source annotation
classes and discards ties. Majority hate maps to $1$. Majority offensive or
normal maps to $0$. This order matters because voting after collapsing the
classes can produce a different label. For hateful examples, available aligned
annotator rationale masks are combined through a positional union. Unicode
NFC normalization and whitespace normalization precede splitting.

\subsection{Splits and duplicate handling}
The implementation uses HateXplain revision
\artifact{01d742279dac941981f53806154481c0e15ee686} and its official split map.
Normalized, case-folded text identifies exact duplicates. With official
assignments, each duplicate group retains the record from the most held-out
split, prioritizing test, then validation, then train. With an unsplit corpus,
duplicate groups with conflicting labels are excluded before seeded splitting.
The resulting HateXplain corpus contains
19,192 posts, distributed as shown in Table~\ref{tab:splits}. Exact matching
does not remove paraphrases, shared templates, or author-level dependencies.

\begin{table}[ht]
\centering
\caption{Deduplicated HateXplain split sizes used by the saved evaluation.}
\label{tab:splits}
\begin{tabular}{lr}
\toprule
Split & Posts\\
\midrule
Train & 15,348\\
Validation & 1,921\\
Test & 1,923\\
\bottomrule
\end{tabular}
\end{table}

\subsection{Frozen classifier used in the prototype}
The application loads the public model
\artifact{pajola/hate-speech-detection-lr-tfidf} at revision
\artifact{8523992895f5145361e49212af7f1d6c94b8f7a9} \cite{checkpoint}.
The artifact is a scikit-learn pipeline containing TF-IDF vectorization and
logistic regression. Its source classes are hate, offensive, and neither.
The wrapper exposes a binary interface by assigning the hate-class probability
to $p_1(x)$ and the sum of the other probabilities to $p_0(x)$. Predictions use
\begin{equation}
\widehat{y}(x)=\mathbb{I}[p_1(x)\geq\tau],\qquad \tau=0.50.
\end{equation}
This decision rule can differ from taking the largest of the original three
class probabilities.

The publisher's preprocessing lowercases text and removes URLs, mentions,
hashtags, the retweet marker, punctuation, and extra whitespace. Those rules
are reproduced inside inference, including when the explanation engine
queries perturbed inputs. A vocabulary check signals when a post yields no
recognized TF-IDF features. Such a score reflects the classifier's intercept
and requires especially careful interpretation.

\subsection{Reproducible supervised baseline}
The project also preserves training code for a word- and character-TF-IDF
linear SVM using scikit-learn \cite{sklearn}. Word features cover unigrams and
bigrams, character features cover lengths two through five, and the SVM uses
balanced class weights. Cross-validation calibration maps margins to class
probabilities. This baseline has its own classroom walkthrough. The numerical
results in this paper belong to the downloaded logistic regression model.

\clearpage
\section{Explanations and the Feedback Framework}
\subsection{Position-based perturbations}
For a post with $m$ whitespace-delimited tokens, let
$z\in\{0,1\}^{m}$ indicate retained token positions. A rendering function
$h_x(z)$ deletes masked tokens while maintaining the surviving token order.
Repeated occurrences receive separate positions and original-text offsets.
The explanation target is $f_c(z)=p_c(h_x(z))$ for selected class $c$.
Explaining the displayed class makes a positive attribution support the
displayed decision and a negative attribution oppose it.

For SHAP, the implementation uses exact evaluation for at most eight tokens.
Longer inputs use a permutation explainer with a budget of
$5(2m+1)$ evaluations. The baseline is a fully deleted post. The additive
interpretation is
\begin{equation}
f_c(\mathbf{1})\approx f_c(\mathbf{0})+\sum_{i=1}^{m}\phi_i.
\end{equation}
Permutation estimates approximate the feature contributions under this
deletion mechanism \cite{shap}.

The LIME route generates 256 binary masks, including the original input, and
fits a local weighted regression through LIME's base implementation
\cite{lime}. Its distance from the original mask is
$d(z)=100(1-\sqrt{m^{-1}\sum_i z_i})$, with kernel
$K(d)=\exp[-\tfrac12(d/25)^2]$. The resulting coefficients measure local effects
under the sampled perturbation neighborhood. Deletions can produce unnatural
language, so their weights should be interpreted in that experimental context.

\subsection{Rationale agreement}
For a human-highlighted set $H$, let $k=|H|$. The explanation set $A$ contains
the positive-scoring positions among the top $k$ attributions for the hate
class. The implementation computes
\begin{equation}
P_r=\frac{|A\cap H|}{|A|},\quad
R_r=\frac{|A\cap H|}{|H|},\quad
F_{1,r}=\frac{2P_rR_r}{P_r+R_r}.
\end{equation}
Zero-support cases use zero precision where appropriate, and empty human masks
are skipped. The overlap score $|A\cap H|/|A\cup H|$ provides an additional
measure. Document scores are averaged across evaluated posts. Because $k$
comes from the human rationale, this protocol measures agreement at a known
annotation length. It does not test autonomous selection of rationale length.

\subsection{Review records and offline preferences}
The Streamlit interface stores a prediction identifier, normalized text,
checkpoint identifier, score, threshold, predicted label, reviewed label, and
timestamp in a local SQLite database. Confirmations receive reward $+1$ and
corrections receive $-1$. Repeated submission of a prediction identifier is
idempotent. Saving a decision creates a review record for later experimentation.

For the transformer extension, let $y^+$ be the reviewed label and
$y^-=1-y^+$. With a frozen reference classifier $p_{\rm ref}$, define
\begin{equation}
\Delta_\theta(x)=
\log\frac{p_\theta(y^+\mid x)}{p_\theta(y^-\mid x)}-
\log\frac{p_{\rm ref}(y^+\mid x)}{p_{\rm ref}(y^-\mid x)}.
\end{equation}
The implemented experimental objective is
$\mathcal{L}_{\rm pref}=-\mathbb{E}\log\sigma(\beta\Delta_\theta(x))$, adapting
the DPO formulation to two class outputs \cite{dpo}. The stored agreement reward
is an audit signal. The reviewed label determines the preference pair.
The tuning path excludes validation and test text matches, keeps the latest
eligible decision per post, and saves a separate candidate checkpoint.
No measured reward-tuning result is available in this study.

\clearpage
\section{Evaluation and Recorded Results}
\subsection{Protocol and scope}
All numerical results below come from the saved report
\artifact{artifacts/pretrained-hatexplain.json}. The report records seed 42,
the deduplicated split assignment, and the frozen threshold 0.50. Its test set
contains 594 hateful and 1,329 non-hateful posts. The pretrained model's
original training split is undocumented, so possible overlap with HateXplain
prevents treating this evaluation as verified unseen-data generalization.
The project's own training partition was unused in this inference experiment.

Precision measures the hateful fraction among flagged posts, while recall
measures the detected fraction of hateful posts. Macro F1 averages the F1
values of the two classes. Table~\ref{tab:metrics} summarizes the measurements.

\begin{table}[ht]
\centering
\caption{Frozen pretrained classifier results, expressed as percentages.}
\label{tab:metrics}
\begin{tabular}{lrr}
\toprule
Metric & Validation & Test\\
\midrule
Accuracy & 64.71 & 62.77\\
Hate precision & 44.26 & 42.36\\
Hate recall & 55.31 & 56.90\\
Hate F1 & 49.18 & 48.56\\
Macro F1 & 61.07 & 59.69\\
False-positive rate & 31.10 & 34.61\\
False-negative rate & 44.69 & 43.10\\
\bottomrule
\end{tabular}
\end{table}

\begin{table}[ht]
\centering
\caption{Test confusion matrix. Rows are gold labels and columns are predictions.}
\label{tab:confusion}
\begin{tabular}{lrr}
\toprule
 & Predicted non-hateful & Predicted hateful\\
\midrule
Gold non-hateful & 869 & 460\\
Gold hateful & 256 & 338\\
\bottomrule
\end{tabular}
\end{table}

The classifier flags 798 test posts, of which 460 are false positives.
It misses 256 of the 594 hateful posts. An always-non-hateful rule would obtain
$1329/1923=69.11\%$ accuracy, exceeding this model's accuracy while having zero
hate recall. This comparison shows why the reported accuracy requires the
class-specific metrics and error counts. The model's operating point leaves
substantial work for a moderator, and the experiment supplies no evidence that
this threshold optimizes any deployment cost.

\subsection{Explanation measurements}
The report samples 30 posts from 594 hateful test posts with nonempty human
rationales using the seeded sampling procedure. SHAP obtains mean rationale
precision 0.624, recall 0.475, F1 0.513, and intersection-over-union 0.415.
These are averages of per-post metrics. Consequently, the harmonic mean of the
reported mean precision and mean recall need not equal the reported mean F1.
The experiment targets the gold hateful class, including when the classifier
predicts non-hateful. It therefore measures evidence alignment for that class
across the sampled hateful posts.

The small sample and union-based human mask limit interpretation. No confidence
interval, LIME rationale score, controlled explanation ablation, or human
usability study was recorded. The integration checks exercised both explainers
and the review workflow, but operational checks do not establish explanation
usefulness or improvements in moderator decisions.

\clearpage
\section{Discussion, Limitations, and Future Experiments}
\subsection{Target-group diagnostics and validity}
Table~\ref{tab:groups} shows selected slices from the test report. Group labels
are reproduced from the dataset and describe annotated targets. They provide
no verified information about the author of a post. Slices can overlap.
False-positive rates use the negative support $N_0$, while false-negative rates
use positive support $N_1$.

\begin{table}[ht]
\centering
\caption{Selected target slices. Error rates are percentages, with denominators shown.}
\label{tab:groups}
\begin{tabular}{lrrrr}
\toprule
Source target label & $N_0$ & $N_1$ & FPR & FNR\\
\midrule
African & 173 & 261 & 45.09 & 27.59\\
Arab & 57 & 103 & 31.58 & 46.60\\
Asian & 42 & 26 & 35.71 & 69.23\\
Caucasian & 117 & 39 & 62.39 & 25.64\\
Homosexual & 159 & 90 & 64.15 & 31.11\\
Islam & 129 & 149 & 28.68 & 57.05\\
Jewish & 76 & 155 & 39.47 & 44.52\\
Women & 258 & 99 & 28.29 & 44.44\\
\bottomrule
\end{tabular}
\end{table}

The variation warrants error inspection. These observational slices mix
different contexts and vocabulary, and several denominators are small.
They cannot independently identify the cause of a disparity. A fuller audit
would combine bootstrap uncertainty estimates, controlled identity substitutions,
and functional tests \cite{sap,hatecheck}. The English evaluation also leaves
Hindi, Hinglish, and other language varieties unmeasured. The publisher's
punctuation and emoji removal may discard useful context, which motivates
testing alternative preprocessing under controlled training conditions.

\subsection{Planned model and feedback comparisons}
The repository contains supervised transformer and binary preference-tuning
source using PyTorch and Hugging Face Transformers \cite{pytorch,transformers}.
A subsequent experiment should compare a newly fitted SVM and an XLM-R
classifier on identical documented splits. Model selection and any threshold
tuning should use validation data. The test partition should be reserved for
the fixed comparison. LoRA offers a possible compute-saving adaptation method
for later experiments \cite{lora}, while the current transformer source uses
ordinary fine-tuning.

Feedback experiments require independently reviewed posts and adjudication of
ambiguous decisions. A useful comparison would evaluate frozen inference,
supervised retraining on corrections, and binary preference tuning using the
same review budget. Each candidate should be checked for class-specific errors,
calibration, and regressions on examples outside the feedback stream. Feedback
collected only from model flags can introduce selection bias. A random sample
of unflagged posts is needed to investigate missed cases.

The experimental image--caption module fuses frozen vision embeddings with
caption TF-IDF features and provides single-modality comparisons. It still
requires a documented meme dataset, independently assigned splits, and measured
results. No multimodal performance is reported here. Explanations for text
alone also leave open how visual evidence should be presented and assessed.

\subsection{Responsible use and conclusion}
Review storage contains user-entered text, so retention, access, and deletion
policies must be defined before use with operational moderation data. Human
decisions can contain disagreement or bias and require quality controls.
Model and dataset documentation should make these conditions explicit
\cite{modelcards,datasheets}.

The implemented prototype connects frozen hate speech classification, signed
word attributions, and persistent moderator decisions. Its measured F1 of
48.56\%, false-positive burden, and limited rationale sample identify concrete
research priorities. The next stage is to establish independent evaluation
data and test whether contextual modeling and reviewed preferences improve
the complete workflow under a consistent definition of hate speech.

\clearpage
\begin{thebibliography}{99}
\small
\bibitem{davidson}
T. Davidson, D. Warmsley, M. Macy, and I. Weber.
Automated Hate Speech Detection and the Problem of Offensive Language.
\emph{Proceedings of ICWSM}, 2017.
\url{https://arxiv.org/abs/1703.04009}.

\bibitem{hatexplain}
B. Mathew, P. Saha, S. M. Yimam, C. Biemann, P. Goyal, and A. Mukherjee.
HateXplain: A Benchmark Dataset for Explainable Hate Speech Detection.
\emph{Proceedings of AAAI}, 2021.
\url{https://arxiv.org/abs/2012.10289}.

\bibitem{hasoc}
T. Mandl et al.
Overview of the HASOC Track at FIRE 2019: Hate Speech and Offensive Content
Identification in Indo-European Languages.
\emph{Proceedings of FIRE}, 2019.
\url{https://doi.org/10.1145/3368567.3368584}.

\bibitem{attention}
A. Vaswani et al. Attention Is All You Need.
\emph{Advances in Neural Information Processing Systems}, 2017.
\url{https://arxiv.org/abs/1706.03762}.

\bibitem{bert}
J. Devlin, M.-W. Chang, K. Lee, and K. Toutanova.
BERT: Pre-training of Deep Bidirectional Transformers for Language Understanding.
\emph{Proceedings of NAACL-HLT}, 2019.
\url{https://aclanthology.org/N19-1423/}.

\bibitem{hatebert}
T. Caselli, V. Basile, J. Mitrovi\'{c}, and M. Granitzer.
HateBERT: Retraining BERT for Abusive Language Detection in English.
\emph{Proceedings of the Workshop on Online Abuse and Harms}, 2021.
\url{https://arxiv.org/abs/2010.12472}.

\bibitem{xlmr}
A. Conneau et al.
Unsupervised Cross-lingual Representation Learning at Scale.
\emph{Proceedings of ACL}, 2020.
\url{https://arxiv.org/abs/1911.02116}.

\bibitem{memes}
D. Kiela, H. Firooz, A. Mohan, V. Goswami, A. Singh, P. Ringshia, and D. Testuggine.
The Hateful Memes Challenge: Detecting Hate Speech in Multimodal Memes.
\emph{Advances in Neural Information Processing Systems}, 2020.
\url{https://arxiv.org/abs/2005.04790}.

\bibitem{lime}
M. T. Ribeiro, S. Singh, and C. Guestrin.
``Why Should I Trust You?'': Explaining the Predictions of Any Classifier.
\emph{Proceedings of ACM SIGKDD}, 2016.
\url{https://arxiv.org/abs/1602.04938}.

\bibitem{shap}
S. M. Lundberg and S.-I. Lee.
A Unified Approach to Interpreting Model Predictions.
\emph{Advances in Neural Information Processing Systems}, 2017.
\url{https://arxiv.org/abs/1705.07874}.

\bibitem{eraser}
J. DeYoung, S. Jain, N. F. Rajani, E. Lehman, C. Xiong, R. Socher, and B. C. Wallace.
ERASER: A Benchmark to Evaluate Rationalized NLP Models.
\emph{Proceedings of ACL}, 2020.
\url{https://aclanthology.org/2020.acl-main.408/}.

\bibitem{hare}
Y. Yang, J. Kim, Y. Kim, N. Ho, J. Thorne, and S.-Y. Yun.
HARE: Explainable Hate Speech Detection with Step-by-Step Reasoning.
\emph{Findings of EMNLP}, 2023.
\url{https://aclanthology.org/2023.findings-emnlp.365/}.

\bibitem{hatecheck}
P. R\"{o}ttger, B. Vidgen, D. Nguyen, Z. Waseem, H. Margetts, and J. Pierrehumbert.
HateCheck: Functional Tests for Hate Speech Detection Models.
\emph{Proceedings of ACL-IJCNLP}, 2021.
\url{https://aclanthology.org/2021.acl-long.4/}.

\bibitem{ouyang}
L. Ouyang et al.
Training Language Models to Follow Instructions with Human Feedback.
\emph{Advances in Neural Information Processing Systems}, 2022.
\url{https://arxiv.org/abs/2203.02155}.

\bibitem{dpo}
R. Rafailov, A. Sharma, E. Mitchell, S. Ermon, C. D. Manning, and C. Finn.
Direct Preference Optimization: Your Language Model Is Secretly a Reward Model.
\emph{Advances in Neural Information Processing Systems}, 2023.
\url{https://arxiv.org/abs/2305.18290}.

\bibitem{sap}
M. Sap, D. Card, S. Gabriel, Y. Choi, and N. A. Smith.
The Risk of Racial Bias in Hate Speech Detection.
\emph{Proceedings of ACL}, 2019.
\url{https://aclanthology.org/P19-1163/}.

\bibitem{modelcards}
M. Mitchell et al. Model Cards for Model Reporting.
\emph{Proceedings of FAT*}, 2019.
\url{https://arxiv.org/abs/1810.03993}.

\bibitem{datasheets}
T. Gebru, J. Morgenstern, B. Vecchione, J. W. Vaughan, H. Wallach,
H. Daum\'{e} III, and K. Crawford.
Datasheets for Datasets.
\emph{Communications of the ACM}, 64(12), 2021.
\url{https://arxiv.org/abs/1803.09010}.

\bibitem{sklearn}
F. Pedregosa et al. Scikit-learn: Machine Learning in Python.
\emph{Journal of Machine Learning Research}, 12, 2011.
\url{https://jmlr.org/papers/v12/pedregosa11a.html}.

\bibitem{pytorch}
A. Paszke et al.
PyTorch: An Imperative Style, High-Performance Deep Learning Library.
\emph{Advances in Neural Information Processing Systems}, 2019.
\url{https://arxiv.org/abs/1912.01703}.

\bibitem{transformers}
T. Wolf et al. Transformers: State-of-the-Art Natural Language Processing.
\emph{Proceedings of EMNLP System Demonstrations}, 2020.
\url{https://aclanthology.org/2020.emnlp-demos.6/}.

\bibitem{lora}
E. J. Hu et al. LoRA: Low-Rank Adaptation of Large Language Models.
\emph{Proceedings of ICLR}, 2022.
\url{https://arxiv.org/abs/2106.09685}.

\bibitem{checkpoint}
pajola. Hate Speech Detection LR-TFIDF.
Hugging Face model card and pretrained artifact, revision
\artifact{8523992895f5145361e49212af7f1d6c94b8f7a9}.
Accessed November 17, 2026.
\url{https://huggingface.co/pajola/hate-speech-detection-lr-tfidf}.
\end{thebibliography}
\end{document}
